% R15 component; only input paths and enumerated location/grammar prose are edited.
% Source: revisions/2026-10-04-r15/manuscript/method_inference_proof.tex; commit 1cb3cc9efe135966ec228dfaf51c6841a6dfc97c.
\subsection{Proof of Proposition~\ref{prop:r15methodinterval}}
\label{sec:r15methodintervalproof}

Condition on the complete executed policy ensemble, including all online
checking and stopping decisions. The final confirmation arrays are independent
of that information. For $C>0$, set
$Z_{sj}=(Y_{sj}+C)/(2C)$. These variables lie in $[0,1]$ and are independent;
their distributions may differ across streams. Apply
\citet[Theorem~11]{maurer2009} to the pooled vector and to its complement,
allocating $\delta/2$ to each sign. After rescaling, the sampling radius is
the first two terms in \eqref{eq:r15methodinterval}. Equal path counts give
\[
 \mathbb E\bar Y=
 \frac{1}{|\mathcal S|}\sum_{s\in\mathcal S}\mathbb E Y_{s1}.
\]
The triangle inequality and the stream-specific transfer accounts bound the
distance between this expression and the desired finite-distribution mean
by $\bar\beta$. This proves the interval claim. When $C=0$, the clipped
observations vanish identically and the transfer account alone suffices.
The conditional bound also holds unconditionally by integration over the
executed ensemble.

At an online checkpoint, condition instead on the preceding execution history.
The policy at that checkpoint is fixed under this conditioning, and its new
checking bank is independent. The same argument applies to that registered
checkpoint event. A union bound over the predeclared events covers the selected
stopping time. Final method inference uses independent confirmation arrays,
so it does not treat the array that triggered a favorable stop as an ordinary
fixed-policy sample. Summing the four family allocations establishes the
stated $0.05$ error bound. On the simultaneous stream event,
\eqref{eq:r15attainment} follows by counting the streams whose lower or upper
endpoints cross the target.\hfill Q.E.D.
